\section{Quản lý Gói Phần mềm (Package Management)}

\subsection{Cập nhật và Cài đặt gói phần mềm}

\textbf{Cách thực hiện:}
Tiến hành cập nhật danh sách các gói phần mềm trên hệ thống với lệnh \texttt{update}, sau đó cài đặt dịch vụ \texttt{vsftpd}.

\begin{lstlisting}[language=bash]
# Cap nhat danh sach
sudo apt update

# Cai dat vsftpd
sudo apt install vsftpd -y
\end{lstlisting}

\textbf{Kết quả:}
Hệ thống lấy danh sách thành công và tự động tải về, cài đặt gói \texttt{vsftpd}.

\begin{figure}[H]
    \centering
    \begin{subfigure}[b]{0.48\textwidth}
        \centering
        \includegraphics[width=\textwidth]{04/apt-update.png}
        \caption{Cập nhật APT}
    \end{subfigure}
    \hfill
    \begin{subfigure}[b]{0.48\textwidth}
        \centering
        \includegraphics[width=\textwidth]{04/apt-install-vsftpd.png}
        \caption{Cài đặt vsftpd}
    \end{subfigure}
    \caption{Cập nhật và cài đặt phần mềm}
\end{figure}

\subsection{Tìm kiếm, Xem chi tiết và Gỡ bỏ}

\textbf{Cách thực hiện:}
- Tìm kiếm từ khóa \texttt{net-tools} trong kho lưu trữ.
- Hiển thị thông tin chi tiết về gói \texttt{vsftpd} vừa cài.
- Thực hiện lệnh xóa hoàn toàn \texttt{vsftpd} khỏi hệ thống.

\begin{lstlisting}[language=bash]
apt search net-tools
apt show vsftpd
sudo apt remove vsftpd -y
\end{lstlisting}

\textbf{Kết quả:}
Các lệnh trả về chính xác kết quả tìm kiếm và thông tin package. Lệnh \texttt{remove} gỡ bỏ thành công.

\begin{figure}[H]
    \centering
    \begin{subfigure}[b]{0.48\textwidth}
        \centering
        \includegraphics[width=\textwidth]{04/apt-search-nettools.png}
        \caption{Tìm kiếm net-tools}
    \end{subfigure}
    \hfill
    \begin{subfigure}[b]{0.48\textwidth}
        \centering
        \includegraphics[width=\textwidth]{04/apt-show-vsftpd.png}
        \caption{Xem chi tiết vsftpd}
    \end{subfigure}
    
    \vspace{0.5cm}
    \begin{subfigure}[b]{0.48\textwidth}
        \centering
        \includegraphics[width=\textwidth]{04/apt-remove-vsftpd.png}
        \caption{Gỡ bỏ vsftpd}
    \end{subfigure}
    \caption{Tra cứu và Gỡ cài đặt}
\end{figure}

\subsection{Quản lý Dependency và Xử lý sự cố}

\textbf{Cách thực hiện:}
- Liệt kê các gói liên quan tới \texttt{python} đã cài trên máy.
- Xem danh sách các gói phụ thuộc (dependencies) của \texttt{vim}.
- Sử dụng \texttt{--fix-broken install} để sửa lỗi package và \texttt{autoremove} để dọn dẹp hệ thống.

\begin{lstlisting}[language=bash]
# Liet ke goi python
dpkg -l | grep python

# Xem phu thuoc cua vim
apt depends vim

# Sua loi va don dep
sudo apt --fix-broken install
sudo apt autoremove -y
\end{lstlisting}

\textbf{Kết quả:}
Quản lý hiệu quả được danh sách các package đã cài và giải phóng không gian ổ cứng an toàn.

\begin{figure}[H]
    \centering
    \begin{subfigure}[b]{0.48\textwidth}
        \centering
        \includegraphics[width=\textwidth]{04/dpkg-grep-python.png}
        \caption{Liệt kê gói cài đặt}
    \end{subfigure}
    \hfill
    \begin{subfigure}[b]{0.48\textwidth}
        \centering
        \includegraphics[width=\textwidth]{04/apt-depends-vim.png}
        \caption{Phụ thuộc của vim}
    \end{subfigure}
    
    \vspace{0.5cm}
    \begin{subfigure}[b]{0.48\textwidth}
        \centering
        \includegraphics[width=\textwidth]{04/apt-fix-broken.png}
        \caption{Sửa lỗi package}
    \end{subfigure}
    \hfill
    \begin{subfigure}[b]{0.48\textwidth}
        \centering
        \includegraphics[width=\textwidth]{04/apt-autoremove.png}
        \caption{Lệnh dọn dẹp autoremove}
    \end{subfigure}
    \caption{Bảo trì hệ thống Package}
\end{figure}
